\documentclass[11pt,a4paper]{article}

\usepackage{fontspec}
\usepackage{polyglossia}
\setmainlanguage{russian}
\setotherlanguage{english}

\setmainfont{DejaVuSerif}[
  Path=fonts/, Extension=.ttf,
  UprightFont=*, BoldFont=*-Bold,
  ItalicFont=*-Italic, BoldItalicFont=*-BoldItalic,
  Scale=0.92]
\setsansfont{DejaVuSans}[
  Path=fonts/, Extension=.ttf,
  UprightFont=*, BoldFont=*-Bold,
  ItalicFont=*-Oblique, BoldItalicFont=*-BoldOblique,
  Scale=0.92]
\setmonofont{DejaVuSansMono}[
  Path=fonts/, Extension=.ttf,
  UprightFont=*, BoldFont=*-Bold,
  ItalicFont=*-Oblique, BoldItalicFont=*-BoldOblique,
  Scale=0.86]

\usepackage[a4paper,left=2cm,right=2cm,top=2.1cm,bottom=2.1cm]{geometry}
\usepackage{amsmath}
\usepackage{amssymb}
\usepackage{graphicx}
\usepackage{float}
\usepackage{booktabs}
\usepackage{longtable}
\usepackage{array}
\usepackage{enumitem}
\usepackage[font=footnotesize,labelfont=bf]{caption}
\usepackage[unicode,hidelinks]{hyperref}

\setlist{nosep,leftmargin=1.4em}
\setlength{\parskip}{0.4em}
\setlength{\parindent}{0pt}
\renewcommand{\arraystretch}{1.15}

% фигуры крупные, и по умолчанию LaTeX пускал бы на страницу лишь одну,
% оставляя половину полосы пустой; ослабляем ограничения на долю плавающих
\renewcommand{\topfraction}{0.95}
\renewcommand{\bottomfraction}{0.95}
\renewcommand{\textfraction}{0.05}
\renewcommand{\floatpagefraction}{0.7}
\setlength{\floatsep}{8pt plus 2pt minus 2pt}
\setlength{\textfloatsep}{10pt plus 2pt minus 2pt}

\newcommand{\term}[1]{\textbf{#1}}
\newcommand{\code}[1]{\texttt{\small #1}}

\begin{document}

\begin{titlepage}
\centering
\vspace*{3cm}
{\sffamily\Huge\bfseries Непрерывная биржа портфельных заявок\par}
\vspace{0.6cm}
{\sffamily\LARGE симулятор мира, агентская экосистема и аналитика прогона\par}
\vspace{2.5cm}
{\large Отчёт по прогону от 17.08.2026\par}
\vspace{1.2cm}
\begin{tabular}{rl}
длина прогона & 12\,000 тиков \\
зерно генератора & 7 \\
агентов в популяции & 64 (выжило 46) \\
время счёта & 64.6 с \\
\end{tabular}
\vfill
{\small Документ и все графики построены автоматически по записи прогона;\\
конфигурация эксперимента задаётся только в \code{SimConfig}.\par}
\end{titlepage}

\tableofcontents
\clearpage

\section*{О чём этот отчёт}
\addcontentsline{toc}{section}{О чём этот отчёт}

Модель состоит из трёх слоёв, и отчёт устроен так же.

\term{Мир} --- две обычные биржи с лимитными стаканами и батч-аукционом,
цены которых связаны общим якорем генерации заявок. Над якорем стоит
экзогенная \emph{справедливая цена}: случайные <<новости>>, которые никто
из участников не наблюдает напрямую.

\term{Непрерывная биржа} (далее CE) --- площадка портфельных заявок,
которая клирится не сведением встречных заявок, а решением линейной задачи
в логарифмах цен. На ней торгуются четыре актива: \(X_1, Y_1\) с первой
площадки и \(X_2, Y_2\) со второй.

\term{Агенты} --- четыре семейства по 16 копий. Трансляторы \(T_1, T_2\)
переносят цену и ликвидность своей лимитной биржи на CE и хеджируют
накопленный инвентарь об домашний стакан; арбитражёры \(A_X, A_Y\)
торгуют только на CE, стягивая цены двойников с разных площадок.
Копии отличаются гиперпараметрами, и эволюционный отбор постепенно
выключает проигрывающих.

Разделы 1 и 2 описывают устройство модели, с формулами и небольшими
примерами <<на пальцах>>. Раздел 3 --- галерея графиков конкретного
прогона. Раздел 4 --- аналитика: что эти графики означают в числах,
включая проверки корректности. Полная конфигурация прогона вынесена
в приложение.


\clearpage
\section{Симулятор мира}

\subsection{Справедливая цена и якорь}

Экзогенная справедливая цена --- это накопленное произведение лог-нормальных
шоков:
\[
  \Phi_t = \Phi_{t-1}\exp(\varepsilon_t),\qquad
  \varepsilon_t \sim \mathcal{N}(0, \sigma_F^2),\qquad
  \sigma_F = 10\ \text{б.п. за тик.}
\]
За весь прогон это даёт блуждание уровня примерно на
\(\sigma_F\sqrt{T} \approx 11.0\%\). Ни один участник \(\Phi_t\) не
видит: шок входит в модель только через якорь генерации фонового потока.

Якорь --- это то, как рынок отслеживает справедливую цену. Он равен
экспоненциальному среднему взвешенных мидов, к которому каждый тик
применяется тот же шок:
\[
  \bar m_t = \frac{\sum_v m_{v,t}\, d_{v,t}}{\sum_v d_{v,t}},\qquad
  a_t = \bigl[(1-\alpha)\,a_{t-1} + \alpha\,\bar m_t\bigr]\exp(\varepsilon_t),
  \qquad \alpha = 1 - 2^{-1/h},\ h = 20 .
\]
Веса \(d_{v,t}\) --- это глубина у мида в полосе
\(\pm4\): тонкая площадка получает малый вес и не может утащить
общий уровень за собой. Обе биржи в одном тике видят \emph{один и тот же}
якорь --- в этом и состоит их связь; никакого прямого арбитража между
лимитными биржами в модели нет.

\begin{figure}[H]
\centering
\includegraphics[width=0.97\textwidth]{figures/ex_coupling.pdf}
\caption[Связь площадок через общий якорь]{\textbf{Связь площадок через общий якорь.} Разовый шок якоря +2\% переносится в цены обеих площадок за несколько десятков тиков; тонкая биржа шумнее и добирается до нового уровня неровно. Справа — вес толстой биржи в якоре: он равен её доле глубины, поэтому тонкая площадка не может утащить общий уровень за собой.}
\label{fig:ex_coupling}
\end{figure}

\subsection{Одна лимитная биржа}

Внутри тика биржа проходит четыре шага: истечение старых заявок, генерация
фонового потока, батч-аукцион, обновление статистики.

\term{Фоновый поток.} Число заявок за тик \(N \sim \text{Пуассон}(\lambda_v)\).
Для каждой заявки \emph{независимо} разыгрываются сторона (монетка 50/50) и
цена \(p \sim \mathcal{N}(a_t, \sigma_p^2)\), округляемая к ценовой сетке с
шагом 0.01. Именно независимость стороны и цены создаёт сделки: часть
покупок оказывается выше части продаж. Каждая заявка живёт 5 тиков.

\begin{figure}[H]
\centering
\includegraphics[width=0.97\textwidth]{figures/ex_book_formation.pdf}
\caption[Откуда берутся стакан, спред и глубина]{\textbf{Откуда берутся стакан, спред и глубина.} Сторона заявки и её цена разыгрываются независимо, поэтому часть покупок оказывается выше части продаж — они и порождают сделки, а неисполненный остаток ложится в стакан (панель б — усреднённый по прогону профиль). Панель в: чем гуще поток, тем уже спред и больше глубина; вертикальные штрихи — интенсивности двух площадок прогона.}
\label{fig:ex_book_formation}
\end{figure}

\term{Батч-аукцион.} Все сделки тика проходят по единой цене \(p^*\),
максимизирующей исполняемый объём:
\[
  p^* = \arg\max_p \min\bigl(D(p),\, S(p)\bigr),\qquad
  D(p) = \!\!\sum_{\text{покупки } p_i \ge p}\!\! q_i,\qquad
  S(p) = \!\!\sum_{\text{продажи } p_i \le p}\!\! q_i .
\]
При нескольких оптимумах берётся цена, ближайшая к предыдущей. На уровне,
где объёма не хватает всем, заявки исполняются pro-rata --- преимущества
от порядка прихода нет. Если пересечения нет, сделок в тике не происходит
и цена последнего аукциона не меняется.

\begin{figure}[H]
\centering
\includegraphics[width=0.97\textwidth]{figures/ex_auction.pdf}
\caption[Батч-аукцион на лимитной бирже]{\textbf{Батч-аукцион на лимитной бирже.} Слева — заявки одного тика, сложенные в кривые спроса и предложения; цена аукциона p* максимизирует исполняемый объём min(спрос, предложение), и все сделки тика проходят по ней одной. Справа — та же механика при сдвинутом якоре генерации: цена аукциона следует за якорем практически один в один, а объём почти не зависит от сдвига.}
\label{fig:ex_auction}
\end{figure}

\term{Наблюдаемые величины.} Мид \(m = (b + a)/2\) (при пустой стороне ---
цена последнего аукциона), спред \(s = a - b\), глубина у мида
\(d = \sum_{|p_i - m| \le 4} q_i\) и волатильность как EWMA
квадратов лог-приростов мида с полураспадом 10 тиков.

\subsection{Две площадки прогона}

Площадки отличаются только густотой потока: у первой
\(\lambda = 10\) заявок за тик, у второй \(\lambda = 3\).
Всё остальное совпадает. Этого достаточно, чтобы получить <<толстую>> и
<<тонкую>> биржу с качественно разной микроструктурой.

\begin{table}[H]
\centering
\caption{Параметры площадок и то, во что они превращаются в прогоне
(средние по всем 12\,000 тикам).}
\small
\begin{tabular}{lrrrrrrrr}
\toprule
& \multicolumn{3}{c}{задано} & \multicolumn{5}{c}{получилось} \\
\cmidrule(lr){2-4}\cmidrule(lr){5-9}
площадка & поток & TTL & \(\sigma_p\) & спред & глубина & вол., б.п.
& объём & \% тиков со сделкой \\
\midrule
биржа 1 & 10.0 & 5 & 3.00 & 0.881 & 17.9 & 66.1 & 2.92 & 99.0 \\
биржа 2 & 3.0 & 5 & 3.00 & 2.162 & 5.3 & 121.7 & 0.91 & 68.3 \\
\bottomrule
\end{tabular}
\end{table}

Тонкая биржа платит за разреженность потока втрое: спред шире, глубина
втрое меньше, а в 12.2\% тиков в стакане вообще нет одной из сторон ---
в такие тики её транслятор молчит, а хедж невозможен.

\subsection{Оценка хеджа: <<бумажное>> исполнение}

Агентам нужен способ сбросить инвентарь об лимитный стакан. В текущей
версии это делается без влияния на рынок: заявка проходит по книге от
лучшей цены вглубь, средняя цена исполнения возвращается агенту, но сам
стакан не меняется. Такое приближение допустимо ровно потому, что
хеджи малы на фоне оборота площадки (панель г рис.~\ref{fig:world_volume}):
за прогон трансляторы отхеджировали 157.5 X1 при суммарном
объёме аукционов 45\,942 штук.


\clearpage
\section{Агентская структура на непрерывной бирже}

\subsection{Непрерывная биржа портфельных заявок}

Состояние биржи --- вектор логарифмов цен \(S \in \mathbb{R}^n\),
\(P = \exp(S)\). Заявка --- это тройка \((w, z, \lambda)\):

\begin{itemize}
\item \(w\) --- веса по стоимости, \(\sum_i w_i = 0\) (заявка
самофинансируема: сколько стоимости вошло, столько и вышло);
\item \(z\) --- курс безразличия, уровень геометрического индекса
\(\Pi(P) = \prod_i P_i^{w_i}\), при котором заявка ничего не торгует;
\item \(\lambda\) --- агрессивность: нотионал за единицу времени на единицу
лог-расхождения, [X1/тик].
\end{itemize}

Сигнал заявки --- лог-расхождение \(f = \ln z - w \cdot S\); заявка торгует
нотионал \(\lambda f\,dt\) (при \(f > 0\) покупает портфель). Условие
клиринга --- нулевой поток стоимости по каждому активу:
\[
  \sum_a \lambda_a f_a\, w_{a,i} = 0 \quad \forall i
  \qquad\Longleftrightarrow\qquad
  A S = b,\quad A = \sum_a \lambda_a w_a w_a^{\mathsf{T}},\quad
  b = \sum_a \lambda_a \ln z_a\, w_a .
\]
Это в точности взвешенная задача наименьших квадратов
\(S^* = \arg\min_S \sum_a \lambda_a (w_a\cdot S - \ln z_a)^2\): цена
клиринга --- \(\lambda\)-взвешенное среднее котировок в логарифмах.
Поскольку \(\sum_i w_{a,i} = 0\) у каждой заявки, \(A\mathbf{1} = 0\), и
решение определено с точностью до общего масштаба цен --- этот произвол
снимается выбором единицы счёта (здесь \(P_{X_1} \equiv 1\)).

Система решается в приращениях от текущих цен, а направления, на которые
никто не выставился, сохраняют прежнюю цену точно --- без регуляризации.
Ранг матрицы \(A\) и есть число ценовых направлений, которые книга
определяет в этот тик (в прогоне обычно 3 из 4).

\begin{figure}[H]
\centering
\includegraphics[width=0.97\textwidth]{figures/ex_ce_clearing.pdf}
\caption[Клиринг непрерывной биржи]{\textbf{Клиринг непрерывной биржи.} Две встречные заявки с ценами безразличия 105 и 95: при равных λ клиринг даёт их геометрическое среднее, при перекосе λ — сдвигается к более агрессивной стороне (левая панель, логарифмическая ось). Справа видно, что цена клиринга — это ровно тот уровень, на котором суммарный поток стоимости обеих заявок обращается в ноль.}
\label{fig:ex_ce_clearing}
\end{figure}

\subsection{Четыре роли}

\begin{table}[H]
\centering
\caption{Портфели заявок. Нога \(+1\) --- покупаемый актив, \(-1\) --- чем
за него платят. У трансляторов индекс портфеля равен цене \(Y\) в кассовом
активе площадки и напрямую сравним с мидом домашней биржи; у арбитражёров
безразличие достигается при курсе 1.}
\small
\begin{tabular}{llll}
\toprule
роль & портфель & курс безразличия \(z\) & что наблюдает \\
\midrule
\(T_1\) & \(+1\,Y_1,\ -1\,X_1\) & мид биржи 1 со сдвигом на позицию
  & стакан биржи 1 \\
\(T_2\) & \(+1\,Y_2,\ -1\,X_2\) & мид биржи 2 со сдвигом на позицию
  & стакан биржи 2 \\
\(A_X\) & \(+1\,X_1,\ -1\,X_2\) & 1 со сдвигом на позицию
  & только цены CE \\
\(A_Y\) & \(+1\,Y_1,\ -1\,Y_2\) & 1 со сдвигом на позицию
  & только цены CE \\
\bottomrule
\end{tabular}
\end{table}

\subsection{Торговая мощность}

Транслятор опирается на домашний стакан: всё, что он наторгует на CE, ему
рано или поздно хеджировать об эту книгу. Качество книги измеряется
величиной с размерностью денег
\[
  H = \frac{d\, m\, \delta}{\max(s,\ \delta)},\qquad \delta = 0.01,
\]
где \(d\,m\) --- нотионал, лежащий у мида, а \(\delta/s\) --- безразмерная
теснота книги (спред в один тик даёт 1). Мощность всего семейства
\(\Lambda = \kappa_T H\) с \(\kappa_T = 1\) делится между копиями
пропорционально капиталу и обратно пропорционально агрессивности:
\[
  \lambda_i = \Lambda \cdot \frac{\max(C_i,0)}{\sum_j \max(C_j,0)}
              \cdot \frac{1}{h_m^{(i)}},\qquad
  C_i = C_0 + \mathrm{PnL}_i,\quad C_0 = 1\,000 .
\]
Арбитражёру опереться не на что, кроме собственного капитала, поэтому у него
\(\lambda_i = \kappa_A \max(C_i,0)/h_m^{(i)}\) с \(\kappa_A = 0.01\).
Мёртвые копии выпадают из суммы, и их доля мощности автоматически перетекает
живым.

\subsection{Риск, котировка и хедж}

Риск-коэффициент агента:
\[
  g = \frac{h_R\,\gamma\,\sigma^2}{\max(C,\ C_{\min})}\quad [1/X_1],
\]
где \(\gamma\) --- общее для системы неприятие риска, \(\sigma^2\) --- EWMA
дисперсии лог-приростов релевантной цены. Богатый агент спокойнее держит
ту же позицию; пол \(C_{\min}\) не даёт разорившемуся делить на ноль.

Котировка сдвигается против собственной позиции (шейдинг):
\[
  z = z_0 \exp(-g\,q),
\]
где \(z_0\) --- базовая оценка курса (мид для транслятора, 1 для
арбитражёра), а \(q\) --- стоимость позиции по торгуемой ноге. Лонг опускает
котировку --- агент охотнее продаёт; сила пружины пропорциональна отклонению
и исчезает при \(q \to 0\).

\begin{figure}[H]
\centering
\includegraphics[width=0.97\textwidth]{figures/ex_shading.pdf}
\caption[Шейдинг котировки против собственной позиции]{\textbf{Шейдинг котировки против собственной позиции.} Цена безразличия z = мид·exp(−g·q) отклоняется от мида тем сильнее, чем больше накопленная позиция и чем выше h\_R (левая панель). Правая панель переводит это в деньги: при цене CE, равной миду, заявка создаёт поток λ·f, направленный на возврат позиции к нулю; наклон задаётся агрессивностью h\_m.}
\label{fig:ex_shading}
\end{figure}

Хедж включается, когда риск последней единицы позиции превысил потерю на её
сбросе:
\[
  \text{держать невыгодно, если } g|q| > c,\quad c = \tfrac12\,\frac{s}{m},
  \qquad\text{сбрасываем излишек } |q| - \frac{c}{g}.
\]
Никаких дополнительных констант здесь нет: и порог, и цель следуют из
сравнения риска и потери в одних единицах, а гистерезис встроен --- сразу
после хеджа риск равен потере, и повторный триггер не сработает, пока
позиция снова не вырастет.

\begin{figure}[H]
\centering
\includegraphics[width=0.97\textwidth]{figures/ex_hedge.pdf}
\caption[Порог хеджа: риск против полуспреда]{\textbf{Порог хеджа: риск против полуспреда.} Держать позицию невыгодно, когда риск последней единицы g·|q| превышает потерю на её сбросе cost = ½·спред/мид; точка равенства и есть порог |q*| = cost/g (левая панель). Справа — игрушечная траектория: без хеджа позиция блуждает, с хеджем удерживается в коридоре, тем более узком, чем выше h\_R.}
\label{fig:ex_hedge}
\end{figure}

Единственная свободная константа \(\gamma\) калибруется по целевому порогу:
если при нейтральном \(h_R = 1\) требовать \(|q^*| = 0.005\,C\), то
\(\gamma = c/(\sigma^2 \cdot 0.005)\). По состоянию рынка после прогрева
это дало \(\gamma = 13\,547\), то есть порог около 5.0 X1 при
стартовом капитале.

\subsection{Гиперпараметры и эволюционный отбор}

Каждая копия несёт свои \(h_m\) (агрессивность: чем меньше, тем больше
\(\lambda\)) и --- у трансляторов --- \(h_R\) (чувствительность к риску: чем
больше, тем раньше и глубже хедж). Сетка \(4\times4\) даёт 16 копий на
семейство, всего 64 агентов.

С тика 6\,000 каждый тик деактивируется не более одного агента:
берётся тот, у кого убыток за последние 2\,000 тиков наибольший, но только
если (а) он убыточен и за всё время, и (б) он не последний живой в своём
семействе. Деактивированный навсегда перестаёт торговать, его прибыль
замораживается.

\subsection{Учёт и точное разложение прибыли}

Инвентарь агента живёт в балансах CE --- это единственный источник правды.
Прибыль считается по рынку: \(\mathrm{PnL} = \sum_i q_i P_i\) в единицах
\(X_1\). Стартовый капитал \(C_0\) на баланс не кладётся, он живёт только
в формулах доли и риска, поэтому прибыль равна ровно стоимости реального
портфеля.

Внутри тика цены CE меняются ровно один раз --- в клиринге. Отсюда точное
разложение изменения прибыли:
\[
  \Delta \mathrm{PnL}_t
  = \underbrace{\sum_i q_{i,t-1}\bigl(P_{i,t} - P_{i,t-1}\bigr)}
    _{\text{переоценка позиции}}
  + \underbrace{\sum_i \Delta q^{\text{сделки}}_{i}\,P_{i,t}}_{= 0}
  + \underbrace{\sum_i \Delta q^{\text{хедж}}_{i}\,P_{i,t}}
    _{\text{результат хеджа}} .
\]
Средний член равен нулю тождественно: \(\sum_i w_i = 0\), поэтому сделка на
CE в момент исполнения не создаёт и не уничтожает стоимость. Значит, весь
PnL --- это переоценка накопленной позиции плюс то, что агент выиграл или
потерял, скидывая инвентарь по встречной стороне книги. Невязка этого
тождества за весь прогон не превысила \(1.8\cdot 10^{-9}\) X1 --- то есть
разложение на рис.~\ref{fig:agents_structure} не приближённое, а точное.


\clearpage
\section{Графики прогона}

Ниже --- полная галерея по прогону длиной 12\,000 тиков (зерно 7).
Прогон делится на три фазы, они размечены фоном на графиках по времени:

\begin{itemize}
\item \term{прогрев}, тики 1--6000: отбор ещё не включён,
живы все 64 копий, гиперпараметры никак не отобраны;
\item \term{отбор}, тики 6000--9800: состав популяции меняется,
доли мощности перетекают выжившим;
\item \term{стационар}, тики 9800--12000: состав зафиксирован,
режим установившийся.
\end{itemize}

Границу стационара выбирает не человек, а данные: последнее выбытие на тике 9800; дальше состав популяции не меняется.

\subsection{Мир}

\begin{figure}[H]
\centering
\includegraphics[width=0.97\textwidth]{figures/world_price.pdf}
\caption[Справедливая цена и цены обычных бирж]{\textbf{Справедливая цена и цены обычных бирж.} Справедливая цена — экзогенный след новостей, общий для обеих площадок; якорь — то, как рынок её отслеживает. Миды обеих бирж идут за якорем, отклоняясь на величину порядка своей микроструктуры: толстая биржа держится плотно, тонкая гуляет заметно шире (панели в, г).}
\label{fig:world_price}
\end{figure}

\begin{figure}[H]
\centering
\includegraphics[width=0.97\textwidth]{figures/world_book.pdf}
\caption[Поведение стаканов обычных бирж]{\textbf{Поведение стаканов обычных бирж.} Тепловые карты (а, б) показывают, где в каждый момент лежит объём относительно мида: у толстой биржи книга плотная и симметричная, у тонкой — рваная, с провалами и односторонними состояниями. Панель в — усреднённый профиль, панель г — спред во времени: именно он входит в качество книги H и в стоимость хеджа.}
\label{fig:world_book}
\end{figure}

\begin{figure}[H]
\centering
\includegraphics[width=0.97\textwidth]{figures/world_volume.pdf}
\caption[Объёмы торгов на обычных биржах]{\textbf{Объёмы торгов на обычных биржах.} Объём аукциона за тик и накопленный оборот (а, б) прямо пропорциональны интенсивности фонового потока. Панель г отвечает на важный для модели вопрос: хеджи трансляторов — это лишь малая доля оборота площадки, поэтому «бумажный» хедж без влияния на рынок остаётся допустимым приближением.}
\label{fig:world_volume}
\end{figure}

\begin{figure}[H]
\centering
\includegraphics[width=0.97\textwidth]{figures/world_liquidity.pdf}
\caption[Ликвидность и глубина рынка]{\textbf{Ликвидность и глубина рынка.} Глубина у мида (а) и качество книги H (б) — это ровно те величины, из которых складывается торговая мощность семейства трансляторов Λ = κ\_T·H. Панель в — цена немедленного исполнения в зависимости от размера: на тонкой бирже она растёт быстрее и глубины хватает не всегда. Панель г сравнивает рыночную волатильность с волатильностью справедливой цены: микроструктурный шум на порядок больше новостного.}
\label{fig:world_liquidity}
\end{figure}

\subsection{Непрерывная биржа}

\begin{figure}[H]
\centering
\includegraphics[width=0.97\textwidth]{figures/ce_price.pdf}
\caption[Цена непрерывной биржи против мира]{\textbf{Цена непрерывной биржи против мира.} Панели а и б — один и тот же срез в двух режимах: в прогреве гиперпараметры ещё не отобраны и цена CE следует за мидом менее уверенно, в стационаре трансляция заметно точнее. Панель в даёт ошибку трансляции во времени, панель г — её распределение по фазам: видно и смещение, и разброс.}
\label{fig:ce_price}
\end{figure}

\begin{figure}[H]
\centering
\includegraphics[width=0.97\textwidth]{figures/ce_basis.pdf}
\caption[Базисы арбитражёров на непрерывной бирже]{\textbf{Базисы арбитражёров на непрерывной бирже.} Арбитражёры удерживают курс кассовых активов двух площадок около паритета (а) и стягивают базисы X1/X2 и Y1/Y2 к нулю (б, в). Панель г показывает, какое расхождение реально видит заявка каждого типа: у арбитражёров оно на порядки меньше, чем у трансляторов, — они работают с почти закрытым базисом.}
\label{fig:ce_basis}
\end{figure}

\begin{figure}[H]
\centering
\includegraphics[width=0.97\textwidth]{figures/ce_clearing.pdf}
\caption[Оборот и качество клиринга непрерывной биржи]{\textbf{Оборот и качество клиринга непрерывной биржи.} Оборот CE и его разбивка по типам (а, б) показывают, кто на самом деле двигает цены. Панель в — наполнение книги: ранг говорит, сколько ценовых направлений книга определяет в каждый тик. Панель г — контрольная: и невязка клиринга, и невязка разложения PnL держатся на уровне машинной точности, то есть стоимость в системе не теряется и не создаётся.}
\label{fig:ce_clearing}
\end{figure}

\subsection{Агенты}

\begin{figure}[H]
\centering
\includegraphics[width=0.97\textwidth]{figures/agents_pnl.pdf}
\caption[PnL агентов и их копий]{\textbf{PnL агентов и их копий.} Панель а — все 64 траектории: до включения отбора живут все копии, дальше проигрывающие замораживаются. Панель в — итог по копиям с отметкой выбывших, панель г сравнивает заработок в прогреве с заработком в стационаре: точки в правом верхнем углу — те, кто был прибыльным в обоих режимах.}
\label{fig:agents_pnl}
\end{figure}

\begin{figure}[H]
\centering
\includegraphics[width=0.97\textwidth]{figures/agents_turnover.pdf}
\caption[Объёмы торгов в разрезе агентов и копий]{\textbf{Объёмы торгов в разрезе агентов и копий.} Панель а — не ошибка рисунка: четыре линии совпадают, потому что условие нулевого потока стоимости связывает нотионалы семейств в кольцо (см. раздел 4), и каждое семейство даёт ровно четверть оборота CE — это же видно на панели г. Панель б показывает, как оборот распределён между копиями (штриховка — выбывшие), панель в — что копии ложатся на закон λ ∝ 1/h\_m.}
\label{fig:agents_turnover}
\end{figure}

\begin{figure}[H]
\centering
\includegraphics[width=0.97\textwidth]{figures/agents_structure.pdf}
\caption[Структурная аналитика PnL агентов]{\textbf{Структурная аналитика PnL агентов.} Разложение точное: PnL = переоценка позиции + результат хеджа (сделки на CE самофинансируемы и в момент исполнения стоимости не создают). Панель а показывает обе компоненты накопленным итогом, панель б — итог по типам, панель в — каждого транслятора в координатах «переоценка / хедж», панель г связывает результат с размером удерживаемой позиции.}
\label{fig:agents_structure}
\end{figure}

\begin{figure}[H]
\centering
\includegraphics[width=0.97\textwidth]{figures/agents_inventory.pdf}
\caption[Инвентарь агентов и работа хеджа]{\textbf{Инвентарь агентов и работа хеджа.} Порог хеджа (штриховая линия на панели а — медиана по блокам, бледная линия — сам ряд) обратно пропорционален дисперсии, поэтому он не постоянен: на всплесках волатильности он проваливается на два порядка и дотягивается до текущей позиции — именно в эти моменты хедж и срабатывает. В остальное время позицию у нуля держит один лишь шейдинг котировки. Инвентарь записан на конец тика, то есть уже после хеджа. Панель в — распределение |позиции| в стационаре, панель г — частота хеджей и их итог по площадкам.}
\label{fig:agents_inventory}
\end{figure}

\begin{figure}[H]
\centering
\includegraphics[width=0.97\textwidth]{figures/agents_hyper.pdf}
\caption[Доли гиперпараметров в живых копиях]{\textbf{Доли гиперпараметров в живых копиях.} Состав популяции во времени: до тика включения отбора доли равны по построению, дальше отбор перераспределяет их. Панели а–в показывают, какие значения h\_m и h\_R выживают, панель г — сколько копий каждого типа осталось.}
\label{fig:agents_hyper}
\end{figure}

\begin{figure}[H]
\centering
\includegraphics[width=0.97\textwidth]{figures/agents_hyper_pnl.pdf}
\caption[Гиперпараметры, результат и выживаемость]{\textbf{Гиперпараметры, результат и выживаемость.} Тепловые карты (а, б) дают итоговый PnL транслятора в каждой точке сетки (h\_m, h\_R); цветовая шкала обрезана по типичному значению, иначе единственный выживший насыщает палитру и структура проигравших не видна — числа в клетках точные. Панель в — средний PnL арбитражёра в разрезе агрессивности, панель г — результат трансляторов к моменту включения отбора: именно на эти числа отбор и реагирует.}
\label{fig:agents_hyper_pnl}
\end{figure}


\clearpage
\section{Аналитика симулятора}

\subsection{Сводка прогона}

За 12\,000 тиков (64.6 с счёта) непрерывная биржа провернула
2\,473 X1 суммарного оборота, в среднем 0.206 X1 за тик
при 56.1 исполнившихся заявках. Трансляторы сделали 1\,901 хеджей
(872 на первой площадке, 1\,029 на второй) на 157.5 X1.
Отбор выключил 18 копий из 64.

\begin{table}[H]
\centering
\caption{Итоги по семействам: выжившие, прибыль и её точное разложение,
оборот на CE и средний размер удерживаемой позиции.}
\small
\begin{tabular}{lrrrrrr}
\toprule
семейство & выжило & PnL, X1 & переоценка & хедж & оборот, X1
& средняя \(|q|\) \\
\midrule
T1 & 11 / 16 & +0.068 & +0.060 & +0.008 & 620.6 & 0.119 \\
T2 & 3 / 16 & −0.366 & −0.667 & +0.301 & 617.6 & 0.109 \\
AX & 16 / 16 & +0.656 & +0.656 & +0.000 & 617.4 & 0.224 \\
AY & 16 / 16 & +0.716 & +0.716 & +0.000 & 617.4 & 0.218 \\
\bottomrule
\end{tabular}
\end{table}

\subsection{Насколько точно CE переносит цену}

\begin{table}[H]
\centering
\caption{По фазам: ошибка трансляции (среднее / стандартное отклонение,
базисные пункты), разброс базисов арбитражёров и оборот CE.}
\small
\begin{tabular}{llrrrrr}
\toprule
фаза & тики & \(Y_1\) к миду 1 & \(Y_2\) к миду 2 &
\(\sigma\) базиса \(X\) & \(\sigma\) базиса \(Y\) & оборот/тик \\
\midrule
прогрев & 1--6000 & +0.6 / 40.5 & +1.4 / 161.2 & 5.66 & 5.66 & 0.216 \\
отбор & 6000--9800 & +1.4 / 38.3 & +2.1 / 154.8 & 5.18 & 5.19 & 0.203 \\
стационар & 9800--12000 & −2.1 / 32.3 & +15.3 / 147.9 & 4.49 & 4.49 & 0.184 \\
\bottomrule
\end{tabular}
\end{table}

Главный результат виден в правой части таблицы. Ошибка трансляции для
толстой площадки падает с 40.5 б.п. в прогреве до 32.3
б.п. в стационаре, для тонкой --- с 161.2 до 147.9 б.п.
Смещения при этом почти нет (−2.1 б.п. в стационаре): CE не
систематически дороже или дешевле лимитной биржи, она просто шумит вокруг
неё, и отбор гиперпараметров этот шум уменьшает.

Базисы арбитражёров держатся на порядок теснее: стандартное отклонение
4.49 б.п. по паре \(X_1/X_2\) и 4.49 б.п. по
\(Y_1/Y_2\). Это ожидаемо: базис --- разница двух цен \emph{внутри} CE, и
чтобы его закрыть, арбитражёру не нужно ходить ни на какой внешний стакан.
То, что оба числа в каждой строке таблицы совпадают, --- не совпадение
и не ошибка; разбор в п.~4.3.

\subsection{Кольцевое тождество оборотов}

На панели а рис.~\ref{fig:agents_turnover} четыре линии совпадают. Это не
артефакт: так устроен клиринг. Обозначим \(V_k\) суммарный нотионал
семейства \(k\) со знаком. Условие нулевого потока стоимости по каждому из
четырёх активов даёт четыре уравнения:
\[
  X_1:\ -V_{T_1} + V_{A_X} = 0,\qquad
  Y_1:\ +V_{T_1} + V_{A_Y} = 0,\qquad
  X_2:\ -V_{T_2} - V_{A_X} = 0,\qquad
  Y_2:\ +V_{T_2} - V_{A_Y} = 0,
\]
откуда \(V_{A_X} = V_{T_1} = -V_{T_2} = -V_{A_Y}\). Все четыре нотионала
равны по модулю, поэтому каждое семейство даёт ровно четверть оборота CE,
как бы ни отличались его \(\lambda\). Отличие возможно лишь в те тики,
когда копии внутри семейства торгуют в разные стороны и валовой оборот
превышает чистый; в прогоне медианное расхождение между семействами
в точности нулевое.

Практическое следствие: <<увеличить долю рынка>> отдельному семейству
нельзя --- можно только перераспределить оборот между копиями внутри
семейства. Именно это и делает отбор.

То же кольцо объясняет и вторую пару совпадающих чисел в таблице~4:
разбросы базисов \(X\) и \(Y\) равны, и это не опечатка. Из
\(V_{A_X} = -V_{A_Y}\) и \(V = \lambda f\) следует
\(\lambda_{A_X} f_{A_X} = -\lambda_{A_Y} f_{A_Y}\); семейства арбитражёров
устроены одинаково и имеют равную суммарную мощность, поэтому
\(f_{A_X} = -f_{A_Y}\) --- базисы оказываются точными зеркалами. Измеренная
на прогоне корреляция равна \(-0.999993\). Экономический смысл: разрыв
между площадками один, и виден он в двух парах активов с разными знаками;
закрыть их независимо нельзя.

\subsection{Кто зарабатывает и на чём}

Итог прогона: арбитражёры +0.656 и +0.716 X1, трансляторы +0.068 и
−0.366 X1. Асимметрия объясняется разложением PnL.

У арбитражёра компонента хеджа тождественно нулевая --- он не выходит на
лимитные биржи. Весь его заработок --- переоценка позиции: он набирает
базис, когда тот расходится, и получает прибыль, когда трансляторы
возвращают цены к паритету. Его позиция ограничена только шейдингом.

У транслятора картина обратная. Он вынужден периодически скидывать
инвентарь по встречной стороне книги, теряя полуспред; на тонкой площадке
спред шире (2.162 против 0.881), и это прямо переносится в результат.
Показательно семейство \(T_2\): переоценка позиции дала −0.667 X1, а
хедж --- +0.301 X1; итоговый −0.366 X1 --- разность двух больших
величин разного знака. Отсюда и разная выживаемость: из 16 копий каждого
семейства трансляторов осталось 11 и 3, а арбитражёров ---
16 и 16 из 16.

Стоит назвать это прямо: в текущей конфигурации отбор выкашивает
трансляторов почти полностью. Это не поломка механики (тождество разложения
PnL выполняется с машинной точностью), а свойство параметров: расплата за
хедж систематически больше, чем премия за перенос цены. Правило отбора
защищает лишь последнюю копию семейства, поэтому трансляция цен на CE
не исчезает --- но держится на одном агенте на площадку.

\subsection{Что отобрала эволюция}

\begin{table}[H]
\centering
\caption{Средний PnL копии в разрезе значений гиперпараметров: у
трансляторов --- за фазу прогрева (именно по ней отбор и принимает решение)
и за весь прогон, у арбитражёров --- за весь прогон.}
\small
\begin{tabular}{lrrrrr}
\toprule
& \multicolumn{3}{c}{трансляторы} & \multicolumn{2}{c}{арбитражёры} \\
\cmidrule(lr){2-4}\cmidrule(lr){5-6}
значение & PnL за прогрев & PnL за прогон & выжило & PnL за прогон & выжило \\
\midrule
\(h_m = 0.55\) & −0.0103 & −0.0183 & 1 / 8 & +0.0649 & 8 / 8 \\
\(h_m = 0.8\) & −0.0076 & −0.0074 & 4 / 8 & +0.0446 & 8 / 8 \\
\(h_m = 1\) & −0.0062 & −0.0072 & 4 / 8 & +0.0357 & 8 / 8 \\
\(h_m = 1.35\) & −0.0048 & −0.0045 & 5 / 8 & +0.0264 & 8 / 8 \\
\(h_R = 0.55\) & −0.0077 & −0.0152 & 4 / 8 & --- & --- \\
\(h_R = 0.8\) & −0.0073 & −0.0063 & 4 / 8 & --- & --- \\
\(h_R = 1\) & −0.0072 & −0.0073 & 4 / 8 & --- & --- \\
\(h_R = 1.35\) & −0.0066 & −0.0086 & 2 / 8 & --- & --- \\
\bottomrule
\end{tabular}
\end{table}

У арбитражёров зависимость от \(h_m\) монотонная и чистая: чем агрессивнее
копия (меньше \(h_m\)), тем больше её \(\lambda\), тем большую долю
семейного оборота она забирает и тем больше зарабатывает. Отбор здесь
никого не тронул --- прибыльны все.

У трансляторов сигнал слабее и сильно зашумлён: разброс PnL между копиями
одного семейства сопоставим с разбросом между значениями гиперпараметра,
поэтому отбор по окну в 2\,000 тиков выключает копии в порядке, который
ближе к случайному, чем к <<по заслугам>>. Панель г
рис.~\ref{fig:agents_pnl} показывает это прямо: связь между прибылью
в прогреве и прибылью в стационаре у трансляторов слабая.

\subsection{Контроль корректности}

Три независимые проверки, все на данных этого прогона:

\begin{enumerate}
\item \term{Клиринг не теряет стоимость.} Максимальная невязка условия
нулевого потока за весь прогон --- \(1.4\cdot 10^{-12}\) X1 при среднем обороте
0.206 X1 за тик, то есть машинный ноль.
\item \term{Разложение PnL точное.} Максимальное расхождение между
изменением прибыли и суммой её компонент --- \(1.8\cdot 10^{-9}\) X1.
\item \term{Единица счёта неподвижна.} Цена \(X_1\) равна 1 по построению
калибра, и все величины отчёта выражены в ней; ранг книги (3 из 4)
показывает, что три ценовых направления определяются заявками, а четвёртое
и есть выбор единицы счёта.
\end{enumerate}

\subsection{Границы применимости}

Что в модели заведомо упрощено и на что это влияет:

\begin{itemize}
\item \term{Хедж не двигает рынок.} Заявка агента проходит по книге, но
стакан остаётся прежним. Приближение опирается на малость хеджей
относительно оборота площадки; при росте \(\kappa_T\) или числа копий его
придётся заменить настоящим исполнением.
\item \term{Заявка живёт один клиринг.} Агенты перевыставляются каждый тик,
поэтому книга CE не накапливает историю --- динамика полностью
определяется текущими наблюдениями.
\item \term{Отбор однонаправлен.} Копии только выключаются, новые не
рождаются и параметры не мутируют. Это отвечает на вопрос <<кто
проигрывает>>, но не на вопрос <<какие параметры оптимальны>>: для второго
нужен механизм воспроизводства.
\item \term{Один прогон --- одна реализация.} Все числа отчёта получены при
зерне 7. Выводы о механике (тождества, разложения, соотношения
площадок) устойчивы, а конкретные PnL и порядок выбытия копий --- нет.
\end{itemize}


\clearpage
\appendix
\section{Конфигурация прогона}

Таблица ниже --- полный дамп \code{SimConfig}. Этого достаточно, чтобы
воспроизвести и прогон, и весь отчёт: при неизменных конфиге и исходниках
симуляции результат берётся из кэша, при любом изменении --- пересчитывается.

\begin{longtable}{p{0.24\textwidth}p{0.20\textwidth}p{0.46\textwidth}}
\toprule
параметр & значение & смысл \\
\midrule
\endhead
\code{h\_m\_values} & 0.55, 0.8, 1, 1.35 & сетка агрессивности \\
\code{h\_r\_values} & 0.55, 0.8, 1, 1.35 & сетка чувствительности к риску \\
\code{arb\_copies} & 4 & копий арбитражёра на значение h_m \\
\code{c0} & 1000.0 & стартовый виртуальный капитал, X1 \\
\code{kappa\_t} & 1.0 & мощность семейства трансляторов, 1/тик \\
\code{kappa\_a} & 0.01 & мощность арбитражёра от капитала, 1/тик \\
\code{gamma} & None & неприятие риска (None -- калибруется) \\
\code{q\_max\_fraction} & 0.005 & целевой хедж-порог как доля капитала \\
\code{capital\_floor\_frac} & 0.01 & пол капитала в риск-коэффициенте \\
\code{shading\_clamp} & 1.0 & предохранитель на |g q| в экспоненте \\
\code{arb\_shading} & True & шейдинг котировки арбитражёров \\
\code{arb\_vol\_half\_life} & 20.0 & полураспад EWMA волатильности базиса \\
\code{total\_steps} & 12000 & длина прогона, тиков \\
\code{evolution\_start} & 6000 & тик включения отбора \\
\code{evolution\_step} & 100 &  \\
\code{window} & 2000 & окно оценки убытка при отборе \\
\code{book\_snapshot\_every} & 25 & период срезов стакана \\
\code{book\_band} & 8.0 & полуширина полосы среза \\
\code{book\_bin} & 0.25 & ширина ценового бина среза \\
\code{report\_dir} & report & каталог отчёта \\
\code{seed} & 7 & зерно генератора \\
\code{warmup} & 200 & прогрев лимитных бирж, тиков \\
\code{tick\_size} & 0.01 & шаг ценовой сетки \\
\code{initial\_price} & 100.0 & стартовая цена \\
\code{depth\_band} & 4.0 & полоса подсчёта глубины у мида \\
\code{anchor\_half\_life} & 20.0 & полураспад EWMA якоря \\
\code{fundamental\_vol} & 0.001 & волатильность справедливой цены за тик \\
\code{progress\_every} & 500 & период печати прогресса \\
\code{venue1} & \multicolumn{2}{p{0.62\textwidth}}{name=1, arrival\_rate=10.0, order\_size=1.0, order\_ttl=5, price\_std=3.0, ewma\_half\_life=10.0} \\
\code{venue2} & \multicolumn{2}{p{0.62\textwidth}}{name=2, arrival\_rate=3.0, order\_size=1.0, order\_ttl=5, price\_std=3.0, ewma\_half\_life=10.0} \\
\bottomrule
\end{longtable}

\end{document}
